% !TeX root = ../数学分析习题（裴礼文）笔记.tex
% 由主文档通过 \input 载入；请编译主文档。
\setcounter{chapter}{0} % 保留原书章号。
\chapter{一元函数极限}
\label{chap:peiliwen-limits}

本章围绕三个问题展开：极限存在的依据是什么，如何计算极限，以及当普通极限不存在时怎样描述变量的渐近行为。前五节以定义、估计和典型题型为主，随后用上、下极限统一讨论振荡现象，最后回到支撑这些论证的实数完备性。微分、积分与级数方法集中放在相应的复习性内容中，使用时说明所需条件。以下以 $\mathbb{N}=\{1,2,\ldots\}$ 表示正整数集；未特别说明时，数列和函数均取实值，“收敛”指收敛到有限实数。

\section{预备}
\label{sec:pei-preliminaries}

\subsection{函数的范围与反函数}
\label{subsec:pei-inverse-functions}

反函数要解决的问题是：\emph{已知函数值，能否唯一找回原来的自变量？}例如 $f(x)=2x+1$，由 $y=2x+1$ 解出 $x=(y-1)/2$，所以反函数为 $f^{-1}(y)=(y-1)/2$。关键在于能否唯一反解，而不只是交换字母 $x,y$。

\begin{definition}
\label{def:pei-inverse}
设 $f:X\to Y$，其中 $X$ 是\emph{定义域}（domain），$Y$ 是\emph{陪域}（codomain），$f(X)=\{f(x):x\in X\}$ 是实际取到的\emph{值域}（range）。若对每个 $y\in f(X)$，方程 $f(x)=y$ 在 $X$ 中恰有一个解，就把这个解记为 $f^{-1}(y)$，称 $f^{-1}:f(X)\to X$ 为 $f$ 的\emph{反函数}（inverse function）。
\end{definition}

\begin{example}
\label{ex:pei-square-inverse-domain}
平方函数 $f(x)=x^2$ 为什么需要限制定义域后才能求反函数？
\end{example}
\begin{solve}
在整个实数轴上，已知 $y=4$ 不能确定原来是 $x=2$ 还是 $x=-2$，所以没有反函数。若只允许 $x\ge0$，则每个 $y\ge0$ 都唯一对应 $x=\sqrt y$，反函数就是平方根函数；若只允许 $x\le0$，反函数则为 $-\sqrt y$。因此，求反函数时必须明确原来的定义域，以及函数实际能取到哪些值。
\end{solve}

\begin{definition}
\label{def:pei-injective-surjective}
方程 $f(x)=y$ 的解 $x\in X$ 称为 $y$ 的原像。对 $f:X\to Y$，若每个 $y\in Y$ 至多有一个原像，称 $f$ 为\emph{单射}（injection）；若每个 $y\in Y$ 至少有一个原像，即 $f(X)=Y$，称 $f$ 为\emph{满射}（surjection）；两者兼有时称为\emph{双射}（bijection）。
\end{definition}

\begin{proposition}[反函数的存在条件]
\label{prop:pei-inverse-existence}
$f:X\to Y$ 在实际值域 $f(X)$ 上有反函数，当且仅当 $f$ 为单射。在整个陪域 $Y$ 上有反函数的充要条件则是 $f$ 为双射。反函数唯一，并满足
\begin{equation}
    f^{-1}(f(x))=x\quad(x\in X),\qquad
    f(f^{-1}(y))=y\quad(y\in f(X)).
    \label{eq:pei-inverse-identities}
\end{equation}
\end{proposition}
\begin{proof}
$y\in f(X)$ 已经保证有解，单射性再保证解唯一，故反函数存在且唯一。反过来，若两个不同的自变量具有同一函数值，就无法唯一反解。若还要求每个 $y\in Y$ 都有解，就需要满射性。式~\eqref{eq:pei-inverse-identities} 分别表示“先求值再找回自变量”和“先找原像再求值”，都回到出发时的数。
\end{proof}

两式还说明原函数与反函数互为反函数，即 $(f^{-1})^{-1}=f$，这里将 $f$ 看作 $X\to f(X)$。\emph[caution]{反函数 $f^{-1}$ 不是倒数 $1/f$}；例如非负半轴上平方函数的反函数是 $\sqrt y$，不是 $1/y^2$。

\begin{proposition}[反函数的两个常用性质]
\label{prop:pei-composite-inverse}
若 $f:X\to Y$、$g:Y\to Z$ 都为双射，则复合函数 $(g\circ f)(x)=g(f(x))$ 也为双射，且
\begin{equation}
    (g\circ f)^{-1}=f^{-1}\circ g^{-1}.
    \label{eq:pei-composite-inverse}
\end{equation}
若实函数 $f:D\to\mathbb R$ 严格单调，则反函数 $f^{-1}:f(D)\to D$ 存在，且与原函数具有相同的增减方向。
\end{proposition}
\begin{proof}
复合时先作用 $f$，再作用 $g$；反解时先撤销 $g$，再撤销 $f$。具体地，对每个 $z\in Z$，先唯一解得 $y=g^{-1}(z)$，再唯一解得 $x=f^{-1}(y)$，所以原来的 $x$ 就是 $f^{-1}(g^{-1}(z))$。

严格单调使不同自变量的函数值不同，故 $f$ 为单射。若 $f$ 递增，$y_1<y_2$ 的两个原像 $x_1,x_2$ 必满足 $x_1<x_2$，否则 $x_1\ge x_2$ 会推出 $y_1\ge y_2$；若 $f$ 递减，则同理由 $x_1\le x_2$ 推出矛盾，故 $x_1>x_2$。这里的“严格”不能删去，例如 $[0,1]$ 上的常值函数就不能唯一反解。
\end{proof}

\begin{example}
\label{ex:pei-rational-inverse}
设 $f(x)=(x-1)/(x+1)$，$x\ge0$，求反函数并写出它的定义域。
\end{example}
\begin{solve}
先把原式写成 $y=1-2/(x+1)$，由 $x\ge0$ 得 $-1\le y<1$。再由 $y(x+1)=x-1$ 解出 $x$，得到
\[
    f^{-1}(y)=\frac{1+y}{1-y},\qquad -1\le y<1.
\]
对这个范围内的每个 $y$，上式都给出唯一的 $x\ge0$，所以原函数的值域确实为 $[-1,1)$，它也正是反函数的定义域。
\end{solve}

\begin{example}
\label{ex:pei-inverse-order}
设 $F(x)=(2x+1)^3$，$x\in\mathbb R$，按复合求逆的顺序求 $F^{-1}$。
\end{example}
\begin{solve}
正向计算先乘 $2$、加 $1$，再立方；反过来，应先开立方，再减 $1$、除以 $2$。由 $y=(2x+1)^3$ 得 $\sqrt[3]{y}=2x+1$，所以
\[
    F^{-1}(y)=\frac{\sqrt[3]{y}-1}{2},\qquad y\in\mathbb R.
\]
开立方和线性反解对每个实数都给出唯一结果，因此 $F:\mathbb R\to\mathbb R$ 为双射。这个例子说明，求复合函数的反函数时，要从最后一步开始逐步还原。
\end{solve}

\subsection{函数的对称性与周期}
\label{subsec:pei-symmetry-period}

研究对称性与周期，都是在寻找函数值之间的联系：把自变量变为 $-x$，可以比较图像的左右两侧；把自变量平移一个固定距离，则可以判断图像是否重复。

\begin{definition}
\label{def:pei-parity}
设 $f$ 的定义域 $D$ 关于原点对称，即 $x\in D$ 时也有 $-x\in D$。若总有 $f(-x)=f(x)$，称 $f$ 为\emph{偶函数}（even function），其图像关于 $y$ 轴对称；若总有 $f(-x)=-f(x)$，称 $f$ 为\emph{奇函数}（odd function），其图像关于原点对称。例如 $x^2$ 是偶函数，$x^3$ 是奇函数。判断奇偶性前，必须先检查定义域是否对称；若奇函数在 $0$ 处有定义，则由 $f(0)=-f(0)$ 得 $f(0)=0$。
\end{definition}

\begin{proposition}[把函数分成偶部与奇部]
\label{prop:pei-parity-decomposition}
定义域关于原点对称的任意函数，都能唯一写成一个偶函数 $H$ 与一个奇函数 $G$ 的和：
\begin{equation}
    f(x)=H(x)+G(x),\qquad
    H(x)=\frac{f(x)+f(-x)}2,\quad
    G(x)=\frac{f(x)-f(-x)}2.
    \label{eq:pei-parity-decomposition}
\end{equation}
\end{proposition}
\begin{proof}
若 $f(x)=H(x)+G(x)$，把 $x$ 换成 $-x$ 就有 $f(-x)=H(x)-G(x)$。两式相加消去奇部，相减消去偶部，便得到上述公式，因此分解只能如此。反过来，按公式定义的 $H,G$ 确实满足 $H(-x)=H(x)$、$G(-x)=-G(x)$，且两者之和为 $f$，所以分解存在且唯一。
\end{proof}

奇函数若有反函数，反函数也仍是奇函数。因为 $y=f(x)$ 时有 $-y=f(-x)$，所以值域也关于原点对称，而且 $f^{-1}(-y)=-x=-f^{-1}(y)$。

\begin{definition}
\label{def:pei-period}
设 $f:\mathbb R\to\mathbb R$。若存在 $T>0$，使每个 $x\in\mathbb R$ 都满足
\[
    f(x+T)=f(x),
\]
则称 $f$ 为\emph{周期函数}（periodic function），$T$ 为它的一个正周期。若所有正周期中有最小者，则称其为最小正周期。
\end{definition}

周期表示图像沿横轴平移 $T$ 后与原图像重合。例如 $\sin x$ 的 $2\pi$、$4\pi$ 都是周期，其中 $2\pi$ 才是最小正周期；常值函数的每个正数都是周期，因而没有最小正周期。若 $f,g$ 有共同的正周期 $T$，它们的和、差、积也以 $T$ 为周期；只分别知道两者各自有周期，还不足以直接这样推断。

知道周期 $T$ 后，只需知道 $[0,T)$ 上的函数值，就能通过平移得到其余部分。若 $f$ 限制在 $[0,T)$ 上时有反函数 $\varphi$，则限制在 $[kT,(k+1)T)$ 上时，反函数为 $y\mapsto\varphi(y)+kT$，其中 $k\in\mathbb Z$，各段对应同一个值域。

\begin{proposition}[由图像对称找周期]
\label{prop:pei-symmetry-period}
设 $f:\mathbb R\to\mathbb R$。若图像关于两条直线 $x=b$、$x=c$（$b<c$）对称，则 $2(c-b)$ 是一个周期；若 $f$ 为奇函数，且图像关于 $x=a\ne0$ 对称，则 $4|a|$ 是一个周期。这些周期不一定最小。
\end{proposition}
\begin{proof}
关于 $x=b$ 对称，就是 $f(2b-x)=f(x)$。先关于 $x=b$ 反射，再关于 $x=c$ 反射，便把 $x$ 变成 $x+2(c-b)$，而函数值不变，即
\[
    f(x+2(c-b))=f(2c-(2b-x))=f(2b-x)=f(x).
\]
第二种情形中，关于 $x=a$ 对称给出 $f(x+2a)=f(-x)=-f(x)$，即平移一次变号，再平移一次恢复原值：$f(x+4a)=f(x)$。反向平移也保持函数值，因此 $4|a|$ 是一个正周期。
\end{proof}

\begin{example}
\label{ex:pei-parity-split}
把 $f(x)=x^3+2x^2-x+1$（$x\in\mathbb R$）写成一个偶函数与一个奇函数的和。
\end{example}
\begin{solve}
先求 $f(-x)=-x^3+2x^2+x+1$，再用式~\eqref{eq:pei-parity-decomposition} 得
\[
    H(x)=2x^2+1,\qquad G(x)=x^3-x.
\]
这里相加后留下偶次幂项，相减后留下奇次幂项；两部分相加正好还原原函数。
\end{solve}

\begin{example}
\label{ex:pei-bounded-second-difference}
设 $f:\mathbb R\to\mathbb R$ 有界，$h>0$，且对每个 $x$ 都有 $f(x+h)-2f(x)+f(x-h)=0$。证明 $h$ 是 $f$ 的周期。
\end{example}
\begin{solve}
把条件改写成 $f(x+h)-f(x)=f(x)-f(x-h)$，可见沿间距为 $h$ 的点依次前进时，相邻函数值的差都相同。固定 $x$，记这个差为 $d=f(x+h)-f(x)$，便有
\[
    f(x+nh)=f(x)+nd,\qquad n=1,2,\ldots.
\]
若 $d\ne0$，右端随 $n$ 无界，与 $f$ 有界矛盾，故 $d=0$，即 $f(x+h)=f(x)$。这对每个 $x$ 都成立，所以 $h$ 为周期。有界性不能删去，例如 $f(x)=x$ 满足差分条件，却没有正周期。
\end{solve}

\subsection{常用不等式与估计}
\label{subsec:pei-basic-inequalities}

做估计，就是用容易控制的量给复杂表达式找上界或下界。例如，要说明误差很小，就应给误差的绝对值找一个小的上界。本节介绍三种常用工具：拆分绝对值、估计对数，以及比较和与积。选择工具时，要同时看清表达式的结构和所需的不等号方向。

\begin{proposition}[三角不等式与反向三角不等式]
\label{prop:pei-triangle-bounds}
对任意实数 $u,v$，\emph{三角不等式}（triangle inequality）及其反向形式（reverse triangle inequality）分别为
\begin{equation}
    |u+v|\le |u|+|v|,\qquad
    \bigl||u|-|v|\bigr|\le |u-v|.
    \label{eq:pei-triangle-bounds}
\end{equation}
\end{proposition}
\begin{proof}
把 $-|u|\le u\le|u|$ 与 $-|v|\le v\le|v|$ 相加，得到 $-(|u|+|v|)\le u+v\le|u|+|v|$，这就是第一式。再由 $u=(u-v)+v$ 得 $|u|-|v|\le|u-v|$；交换 $u,v$ 得 $|v|-|u|\le|u-v|$，合起来便是第二式。
\end{proof}

第一式把“控制一个和”转成“分别控制各项”，多个加数也可以反复使用。例如，给定 $\delta>0$，若 $|x-a|\le\delta$、$|y-b|\le\delta$，则 $|(x+y)-(a+b)|\le2\delta$。第二式则说明，\emph{取绝对值不会扩大两个数之间的距离}：若 $|x-a|\le\delta$，就也有 $\bigl||x|-|a|\bigr|\le\delta$。

\begin{proposition}[对数的上下界]
\label{prop:pei-log-bounds}
用 $\log$ 表示自然对数。对 $t>-1$，有
\begin{equation}
    \frac{t}{1+t}\le\log(1+t)\le t,
    \label{eq:pei-log-bounds}
\end{equation}
两侧等号均当且仅当 $t=0$ 成立。条件 $t>-1$ 保证对数有定义，也保证分母为正。
\end{proposition}
\begin{proof}
这里用导数作复习性证明。令 $L(s)=s-1-\log s$（$s>0$），则 $L'(s)=(s-1)/s$。所以 $L$ 在 $(0,1)$ 上递减，在 $(1,+\infty)$ 上递增，最小值为 $L(1)=0$，即 $\log s\le s-1$，等号仅在 $s=1$ 时成立。取 $s=1+t$ 得到上界；再取 $s=1/(1+t)$，得到 $-\log(1+t)\le-t/(1+t)$，两边乘以 $-1$ 并将不等号反向，便得下界。
\end{proof}

这组不等式把对数夹在两个代数式之间。例如，对正整数 $k$ 取 $t=1/k$，便有 $1/(k+1)<\log(1+1/k)<1/k$，可用于比较对数项与倒数项。\emph[caution]{除以负数时不等号要反向}，因此后续若要除以 $t$，必须区分 $t>0$ 与 $t<0$。

\begin{theorem}[算术平均与几何平均不等式]
\label{thm:pei-amgm}
设 $n$ 为正整数，$a_1,\ldots,a_n\ge0$，则
\begin{equation}
    \underbrace{(a_1\cdots a_n)^{1/n}}_{\text{几何平均}}
    \le\underbrace{\frac{a_1+\cdots+a_n}{n}}_{\text{算术平均}},
    \label{eq:pei-amgm}
\end{equation}
等号成立当且仅当 $a_1=\cdots=a_n$。它也称\emph{均值不等式}（arithmetic--geometric mean inequality）：当这些非负数的和固定时，它们的乘积在各数相等时达到最大值。
\end{theorem}
\begin{proof}
二元情形来自 $(\sqrt{a_1}-\sqrt{a_2})^2\ge0$。一般情形中，若有某个 $a_i=0$，左端为零，结论成立，且等号只能在所有数均为零时成立。以下设各数都为正，记算术平均为 $A=(a_1+\cdots+a_n)/n$。由式~\eqref{eq:pei-log-bounds} 的上界，取 $t=a_i/A-1>-1$，再对 $i$ 求和，得到
\[
    \log\frac{a_1\cdots a_n}{A^n}
    =\sum_{i=1}^n\log\frac{a_i}{A}
    \le\sum_{i=1}^n\left(\frac{a_i}{A}-1\right)=0.
\]
因此 $a_1\cdots a_n\le A^n$，开 $n$ 次方即得结论。等号要求每个 $a_i/A=1$，也就是各数相等。
\end{proof}

\begin{example}
\label{ex:pei-amgm-minimum}
设 $x>0$，求 $x+4/x$ 的最小值。
\end{example}
\begin{solve}
两个正数 $x$ 与 $4/x$ 的乘积固定为 $4$，因此用二元均值不等式有
\[
    x+\frac4x\ge2\sqrt{x\cdot\frac4x}=4.
\]
等号要求 $x=4/x$，结合 $x>0$ 得 $x=2$。所以最小值为 $4$，在 $x=2$ 时取得；仅有下界还不够，检查等号能够取到，才确定了最小值。
\end{solve}

\begin{example}
\label{ex:pei-log-approximation-error}
设 $|t|\le1/2$，估计用 $t$ 代替 $\log(1+t)$ 时的误差。
\end{example}
\begin{solve}
由式~\eqref{eq:pei-log-bounds}，$\log(1+t)\le t$，所以误差的绝对值等于 $t-\log(1+t)$。同时，$1+t\ge1/2$，故
\[
    0\le t-\log(1+t)
    \le t-\frac{t}{1+t}
    =\frac{t^2}{1+t}\le2t^2.
\]
因此误差不超过 $2t^2$；例如用 $0.01$ 近似 $\log(1.01)$，误差至多为 $0.0002$。这个估计对正、负的 $t$ 都成立，并且给出了随 $t$ 缩小而缩小的误差界。
\end{solve}

\subsection{Wallis 积分与乘积}

以下公式使用定积分与分部积分，可在学习积分后回看。非负整数的\emph{双阶乘}（double factorial）定义为 $n!!=n(n-2)(n-4)\cdots$，并约定 $0!!=(-1)!!=1$。令
\[
    I_n=\int_0^{\pi/2}\sin^n x\,dx
       =\int_0^{\pi/2}\cos^n x\,dx.
\]
变量替换 $x\mapsto\pi/2-x$ 给出两积分相等。对 $n\ge2$ 分部积分，得到 $I_n=(n-1)(I_{n-2}-I_n)$，所以
\begin{equation}
    I_n=\frac{n-1}{n}I_{n-2}
    =\frac{(n-1)!!}{n!!}
    \begin{cases}
        \pi/2,&n\text{ 为偶数},\\
        1,&n\text{ 为奇数}.
    \end{cases}
    \label{eq:pei-wallis-integral}
\end{equation}
这里的初值是 $I_0=\pi/2$、$I_1=1$。由于 $0\le\sin x\le1$，有 $I_{2n+1}\le I_{2n}\le I_{2n-1}$，因而
\[
    1\le\frac{I_{2n}}{I_{2n+1}}
      \le\frac{I_{2n-1}}{I_{2n+1}}=\frac{2n+1}{2n}\longrightarrow1.
\]
代入式~\eqref{eq:pei-wallis-integral}，得到\emph{Wallis 乘积公式}
\begin{equation}
    \lim_{n\to\infty}
    \prod_{k=1}^n\frac{(2k)^2}{(2k-1)(2k+1)}=\frac\pi2.
    \label{eq:pei-wallis-product}
\end{equation}

对非负整数 $m,n$，还可写出统一形式
\begin{equation}
    \int_0^{\pi/2}\sin^m x\cos^n x\,dx
    =\frac{(m-1)!!(n-1)!!}{(m+n)!!}
    \begin{cases}
        \pi/2,&m,n\text{ 均为偶数},\\
        1,&\text{其余情形}.
    \end{cases}
    \label{eq:pei-mixed-wallis}
\end{equation}
记左端为 $J_{m,n}$，当 $n\ge2$ 时，对 $\sin^m x\cos x\cdot\cos^{n-1}x$ 分部积分得 $J_{m,n}=(n-1)J_{m+2,n-2}/(m+1)$。反复降阶至 $J_{m,0}=I_m$ 或 $J_{m,1}=1/(m+1)$，即可得到式~\eqref{eq:pei-mixed-wallis}。

\subsection{确界与单调映射的一个应用}

非空实数集 $E$ 的\emph{上确界}（supremum）$s=\sup E$ 指最小上界，其可操作的刻画为：每个 $x\in E$ 都有 $x\le s$，且每个 $\varepsilon>0$ 都能找到 $x_\varepsilon\in E$ 使 $s-\varepsilon<x_\varepsilon\le s$。\emph{下确界}（infimum）$i=\inf E$ 相应满足 $i\le x$，且任意 $\varepsilon>0$ 都有 $x_\varepsilon\in E$ 使 $i\le x_\varepsilon<i+\varepsilon$。确界不要求属于集合；属于集合时才分别是最大值与最小值。有界非空实数集的确界存在性将在第~\ref{sec:pei-real-completeness} 节集中讨论。

若 $\{x_n\}$、$\{y_n\}$ 有界，则逐项估计 $x_n+y_n\le\sup_kx_k+\sup_ky_k$ 给出 $\sup_n(x_n+y_n)\le\sup_nx_n+\sup_ny_n$；这里一般不能写等号，因为两列的较大值未必出现在同一指标上。若 $A,B$ 是非空有界的非负数集，允许 $x\in A$ 与 $y\in B$ 独立选取，则有 $\sup\{xy:x\in A,y\in B\}=(\sup A)(\sup B)$：一侧由上界性质得到，另一侧由分别逼近两个上确界得到。

\begin{example}
\label{ex:pei-monotone-fixed-point}
设 $f:[a,b]\to[a,b]$ 单调不减，证明存在 $c\in[a,b]$ 使 $f(c)=c$，不要求 $f$ 连续。
\end{example}
\begin{solve}
令 $S=\{x\in[a,b]:x\le f(x)\}$。因 $a\le f(a)$，集合 $S$ 非空且有上界，故 $c=\sup S\in[a,b]$。对任意 $x\in S$，单调性给出 $x\le f(x)\le f(c)$，所以 $c\le f(c)$。再次使用单调性，$f(c)\le f(f(c))$，故 $f(c)\in S$，从而 $f(c)\le c$，合起来得 $f(c)=c$。
\end{solve}

% 第 1.2 节：PDF 第 9--24 页。
\section{用定义证明极限的存在性}
\label{sec:pei-limit-existence}

\subsection{定义、量词与误差估计}

\begin{definition}
\label{def:pei-sequence-limit}
称数列 $\{x_n\}$ 收敛到 $A\in\mathbb{R}$，记为 $x_n\to A$，是指
\[
    \forall\varepsilon>0,\quad\exists N\in\mathbb{N},\quad
    \forall n>N,\quad |x_n-A|<\varepsilon.
\]
设 $a$ 是 $E\subseteq\mathbb{R}$ 的聚点，$f:E\to\mathbb{R}$。称 $x\to a$、$x\in E$ 时 $f(x)\to A$，是指
\[
    \forall\varepsilon>0,\quad\exists\delta>0,\quad
    x\in E,\ 0<|x-a|<\delta\ \Longrightarrow\ |f(x)-A|<\varepsilon.
\]
这里“聚点”意指任意去心邻域中都有 $E$ 的点，详见第~\ref{sec:pei-real-completeness} 节。函数在 $a$ 处的取值不影响此极限。
\end{definition}

还约定 $x_n\to+\infty$ 表示对每个 $M>0$，存在 $N$ 使 $n>N$ 时 $x_n>M$；$x_n\to-\infty$ 表示 $-x_n\to+\infty$。若 $|x_n|\to+\infty$，称它为\emph{无穷大量}，不要求各项最终同号。函数趋于 $\pm\infty$ 时，把上述尾部条件换成相应的去心邻域条件；讨论 $x\to+\infty$ 的函数极限时，则把 $0<|x-a|<\delta$ 换成 $x>R$。这些广义极限与有限收敛应加以区分。

证明的核心是把误差控制在事先给定的 $\varepsilon$ 内。若可以解出 $|x_n-A|<\varepsilon$ 对 $n$ 的充分条件，就直接选取 $N$；若直接求解困难，则先放大为简单的上界；还可先要求 $n>N_0$ 或 $|x-a|<\delta_0$，取得局部简化后再分配误差。\emph[tip]{所选阈值只需保证结论，不要求最小或最优。}例如，$N$ 可以依赖 $\varepsilon$，但一经选定，就必须同时控制所有 $n>N$。

\begin{example}
\label{ex:pei-root-n-definition}
用定义证明 $n^{1/n}\to1$。
\end{example}
\begin{solve}
令 $u_n=n^{1/n}-1\ge0$。对 $n\ge2$，二项式展开中各项非负，故
\[
    n=(1+u_n)^n\ge\frac{n(n-1)}2u_n^2,
    \qquad 0\le u_n\le\sqrt{\frac2{n-1}}.
\]
给定 $\varepsilon>0$，取整数 $N>1+2/\varepsilon^2$，则 $n>N$ 时 $|n^{1/n}-1|<\varepsilon$。这里先引入误差变量，再通过二项式中的一个正项取得可解的上界。
\end{solve}

\begin{example}
\label{ex:pei-delta-stages}
用定义证明 $\displaystyle\lim_{x\to1}\frac7{16x^2-9}=1$。
\end{example}
\begin{solve}
先限制 $|x-1|<1/8$，于是 $7/8<x<9/8$，从而 $16x^2-9>3$、$|x+1|<3$。因此
\[
    \left|\frac7{16x^2-9}-1\right|
    =\frac{16|x-1||x+1|}{|16x^2-9|}<16|x-1|.
\]
给定 $\varepsilon>0$，取 $\delta=\min\{1/8,\varepsilon/16\}$ 即可。先隔开分母的零点，再控制误差，是分步估计的典型做法。
\end{solve}

不收敛到某个给定实数 $A$ 的否定形式是
\[
    \exists\varepsilon_0>0,\quad\forall N\in\mathbb{N},\quad
    \exists n>N,\quad |x_n-A|\ge\varepsilon_0.
\]
它只排除这个 $A$；要证明数列没有有限极限，还须排除所有实数候选值，或改用 Cauchy 准则的否定形式。量词取反时，依次交换“任意”与“存在”，同时否定最后的不等式，不能只把 $<$ 换成 $\ge$。

\subsection{Cauchy 准则与单调有界原理}

\begin{theorem}[Cauchy 收敛准则]
\label{thm:pei-cauchy}
实数列 $\{x_n\}$ 收敛，当且仅当
\begin{equation}
    \forall\varepsilon>0,\quad\exists N\in\mathbb{N},\quad
    \forall n>N,\ \forall p\in\mathbb{N},\quad
    |x_{n+p}-x_n|<\varepsilon.
    \label{eq:pei-cauchy}
\end{equation}
满足该条件的数列称为\emph{Cauchy 列}（Cauchy sequence），也称基本列。
\end{theorem}

必要性由 $|x_{n+p}-x_n|\le|x_{n+p}-A|+|x_n-A|$ 得到；充分性依赖实数的完备性，证明见第~\ref{sec:pei-real-completeness} 节。其优点是不需要预先猜出 $A$。\emph[caution]{式~\eqref{eq:pei-cauchy} 中的 $N$ 必须与 $p$ 无关。}对每个固定的 $p$ 有 $x_{n+p}-x_n\to0$ 只是必要条件，例如 $x_n=\sqrt n$ 满足这一条件，却趋于 $+\infty$。

\begin{example}
\label{ex:pei-harmonic-divergence}
证明调和部分和 $H_n=\sum_{k=1}^n1/k$ 不收敛。
\end{example}
\begin{solve}
对每个 $n\ge1$，有 $H_{2n}-H_n=\sum_{k=n+1}^{2n}1/k\ge1/2$。取 $\varepsilon_0=1/2$，任意阈值之后总能选取这样的两项，故不满足 Cauchy 准则。这里必须允许步长 $p=n$ 随 $n$ 改变；固定 $p$ 时差值反而趋于零。
\end{solve}

\begin{theorem}[单调有界原理]
\label{thm:pei-monotone}
实数列单调不减且有上界时，收敛到其全体项的上确界；单调不增且有下界时，收敛到其全体项的下确界。
\end{theorem}
\begin{proof}
只证递增情形。令 $A=\sup\{x_n:n\in\mathbb{N}\}$。给定 $\varepsilon>0$，由确界刻画存在 $N$ 使 $A-\varepsilon<x_N\le A$。对所有 $n>N$，单调性给出 $A-\varepsilon<x_N\le x_n\le A$，即 $|x_n-A|<\varepsilon$。递减情形对 $-x_n$ 应用同一结论即可。
\end{proof}

\begin{example}
\label{ex:pei-euler-constant}
证明 $H_n-\log n$ 收敛。
\end{example}
\begin{solve}
由式~\eqref{eq:pei-log-bounds}，
\[
    (H_{n+1}-\log(n+1))-(H_n-\log n)
    =\frac1{n+1}-\log\left(1+\frac1n\right)<0.
\]
又有 $\log n=\sum_{k=1}^{n-1}\log(1+1/k)<H_{n-1}<H_n$，故该数列有下界零，因而收敛。记其极限为 $\gamma_{\mathrm E}$，称为\emph{Euler 常数}，于是 $H_n=\log n+\gamma_{\mathrm E}+o(1)$；记号 $o(1)$ 表示趋于零的量。此结论只给出存在性，没有把未知极限当作已知数值。
\end{solve}

\subsection{子列、重新编号与函数极限}

\begin{theorem}[子列判别与 Heine 判别]
\label{thm:pei-heine}
数列 $x_n\to A$ 当且仅当每个子列 $x_{n_k}$ 都趋于同一个 $A$，其中 $n_1<n_2<\cdots$。函数 $f:E\to\mathbb{R}$ 在聚点 $a$ 处极限为 $A$，当且仅当对每个满足 $t_n\in E\setminus\{a\}$、$t_n\to a$ 的数列，都有 $f(t_n)\to A$。
\end{theorem}
\begin{proof}
从原极限到子列极限，直接使用定义即可。函数情形的逆向证明采用反证法：若极限不是 $A$，存在 $\varepsilon_0>0$，使每个去心邻域内都有 $t$ 满足 $|f(t)-A|\ge\varepsilon_0$。逐次在半径 $1/n$ 的邻域中选取 $t_n$，便有 $t_n\to a$ 而 $f(t_n)$ 不趋于 $A$。数列情形用不断增大的指标作同样选择。
\end{proof}

有限个子列若覆盖原数列的全部指标，且它们均趋于同一个 $A$，则原数列趋于 $A$；证明时把各子列的阈值换算为原指标后取最大值。特别，$x_n\to A$ 等价于 $x_{2n}\to A$ 且 $x_{2n+1}\to A$。仅有每个子列都能再选出某个收敛子列，并不能保证原数列收敛，例如 $(-1)^n$；若所有这些再选出的子列都趋于同一个指定的 $A$，则可用反证法得到原列收敛到 $A$。

收敛数列还允许适当重新编号。若 $g:\mathbb{N}\to\mathbb{N}$ 的每个纤维 $g^{-1}(\{k\})$ 都是有限集，则 $g(m)\to+\infty$：对每个 $K$，集合 $\{m:g(m)\le K\}$ 是有限个有限集的并。因此 $a_n\to A$ 蕴含 $a_{g(m)}\to A$。重新编号后的指标不必单调，故它未必是通常定义下的子列。

\subsection{平均、拟合与有限运算}

\begin{theorem}[加权平均保持极限]
\label{thm:pei-weighted-cesaro}
设 $a_n\to a$，$p_n\ge0$，$P_n=\sum_{k=1}^np_k\to+\infty$。则
\begin{equation}
    \frac{\sum_{k=1}^np_ka_k}{P_n}\longrightarrow a.
    \label{eq:pei-weighted-cesaro}
\end{equation}
其中只需在 $P_n>0$ 的充分大指标上讨论。
\end{theorem}
\begin{proof}
给定 $\varepsilon>0$，先取 $K$ 使 $k>K$ 时 $|a_k-a|<\varepsilon/2$。分开有限首段与尾段，有
\[
    \left|\frac{\sum_{k=1}^np_ka_k}{P_n}-a\right|
    \le\frac{\sum_{k=1}^Kp_k|a_k-a|}{P_n}
       +\frac\varepsilon2\frac{\sum_{k=K+1}^np_k}{P_n}.
\]
固定 $K$ 后，第一项趋于零，第二项不超过 $\varepsilon/2$，故结论成立。
\end{proof}

取 $p_k=1$ 就是算术平均的\emph{Cesàro 定理}，取 $p_k=k$ 得到 $\sum_{k=1}^nka_k/\sum_{k=1}^nk\to a$。更一般，若三角形系数 $c_{n,k}$ 满足 $\sum_{k=1}^nc_{n,k}=1$、$\sup_n\sum_{k=1}^n|c_{n,k}|<+\infty$，且每个固定 $k$ 有 $c_{n,k}\to0$，则 $\sum_{k=1}^nc_{n,k}a_k\to a$。证明仍是上述分段估计：有限首段由逐列趋零控制，尾段由绝对值和的一致上界控制；这也说明为什么带符号权重需要额外限制。

\begin{example}
\label{ex:pei-fitting-sine}
设 $f(t)/t\to1$（$t\to0$，$t\ne0$），固定 $a\ne0$。证明
\[
    \sum_{k=1}^nf\left(\frac{(2k-1)a}{n^2}\right)\longrightarrow a.
\]
\end{example}
\begin{solve}
记 $t_{n,k}=(2k-1)a/n^2$，则 $\sum_kt_{n,k}=a$，且 $\max_k|t_{n,k}|\le2|a|/n\to0$。给定 $\eta>0$，充分大 $n$ 时所有 $k\le n$ 同时满足 $|f(t_{n,k})-t_{n,k}|\le\eta|t_{n,k}|$，于是总误差不超过 $\eta\sum_k|t_{n,k}|=\eta|a|$。因 $\eta$ 任意，结论成立。此法把目标值 $a$ 写成与原和式对应的形式；取 $f(t)=\sin t$ 就得到典型的三角和极限。
\end{solve}

有限极限具有和、差、积运算法则；商法则还要求分母极限不为零。若 $x_n>0$、$x_n\to a>0$、$y_n\to b\in\mathbb{R}$，则由 $x_n^{y_n}=\exp(y_n\log x_n)$ 得 $x_n^{y_n}\to a^b$。\emph[caution]{这些法则不能直接推广到项数随 $n$ 增加的和或积。}例如令 $u_{n,k}=1$（$k=n$）、$u_{n,k}=0$（$k\ne n$），每个固定 $k$ 都有 $u_{n,k}\to0$，但 $\sum_{k=1}^nu_{n,k}=1$；例~\ref{ex:pei-fitting-sine} 中对整行误差的控制正是不可省去的步骤。

% 第 1.3 节：PDF 第 24--46 页。
\section{求极限值的若干方法}
\label{sec:pei-computing-limits}

\subsection{等价替换、相消与初等变形}

当 $v$ 在极限过程的充分后期不为零，且 $u/v\to1$，称 $u$ 与 $v$ \emph{渐近等价}（asymptotically equivalent），记为 $u\sim v$；$r=o(v)$ 表示 $r/v\to0$，$r=O(v)$ 表示该比值在充分后期有界。下列等价式均针对 $t\to0$：
\begin{equation}
    \begin{gathered}
        \sin t\sim\tan t\sim\arcsin t\sim\arctan t\sim t,\\
        \log(1+t)\sim e^t-1\sim t,\qquad
        1-\cos t\sim\tfrac12t^2,\\
        a^t-1\sim t\log a\quad(a>0,\ a\ne1),\qquad
        (1+t)^\alpha-1\sim\alpha t\quad(\alpha\ne0).
    \end{gathered}
    \label{eq:pei-equivalents}
\end{equation}
其中三角函数式可由基本三角极限得到，其余可由相应函数在零点的导数定义得到；使用反三角函数时取零点附近的通常分支。

\begin{proposition}[等价替换的范围]
\label{prop:pei-equivalent-replacement}
若 $u\sim v$，且相关商有定义，则在乘除结构中把因子 $u$ 换成 $v$，不改变存在的极限。但在 $u-w$ 中作同样替换，需要进一步验证 $u-v=o(u-w)$，或给出足以控制相消后误差的其他估计。
\end{proposition}

例如 $\displaystyle\lim_{t\to0}\frac{\sin t-t}{t^3}=-1/6$，若仅把 $\sin t$ 换成 $t$，会错误地得到零；同样，$\displaystyle\lim_{t\to0}\frac{1-\cos t-t^2/2}{t^4}=-1/24$，主项抵消后必须保留四阶信息。这里的数值可由本节后面的 Taylor 公式核查。相反，$\log(1+\sin t)/t$ 可分解为 $[\log(1+\sin t)/\sin t](\sin t/t)$，两个因子都趋于一，所以原极限为一。

先作精确恒等变形，往往比近似更直接。对 $n\ge2$，
\[
    \prod_{k=2}^n\left(1-\frac1{k^2}\right)
    =\prod_{k=2}^n\frac{(k-1)(k+1)}{k^2}
    =\frac{n+1}{2n}\longrightarrow\frac12.
\]
由倍角公式，对固定 $x\ne0$ 和充分大 $n$，
\[
    \prod_{k=1}^n\cos\frac{x}{2^k}
    =\frac{\sin x}{2^n\sin(x/2^n)}\longrightarrow\frac{\sin x}{x};
\]
当 $x=0$ 时各乘积为一。根式差通常先乘共轭式，例如 $\sqrt{n^2+n}-n=1/(\sqrt{1+1/n}+1)\to1/2$；由此再利用三角函数的连续性可得 $\sin^2(\pi\sqrt{n^2+n})\to1$。

\subsection{已知极限、变量替换与夹逼}

\begin{proposition}[幂式极限]
\label{prop:pei-power-limit}
若 $u\to0$、$1+u>0$，且 $uv\to L\in\mathbb{R}$，则 $(1+u)^v\to e^L$。更一般，对任意正底数 $F$，都应先研究 $G\log F$；若它趋于 $L$，则 $F^G\to e^L$。
\end{proposition}
\begin{proof}
对 $u\ne0$，有 $v\log(1+u)=uv\,[\log(1+u)/u]\to L$，在 $u=0$ 的位置左端为零且与原式相容；再利用指数函数连续性即可。若 $u$ 在无穷多位置为零，则 $uv$ 在这些位置也为零，所以有限极限 $L$ 必为零，论证仍成立。
\end{proof}

例如 $a,b>0$ 时，$\left[(a^{1/n}+b^{1/n})/2\right]^n\to\sqrt{ab}$，因为底数为 $1+(\log a+\log b)/(2n)+o(1/n)$。若 $a=0<b$，原式为 $b/2^n\to0$，与 $\sqrt{ab}=0$ 相符；不能在这一边界情形使用 $\log a$。这说明参数边界需要单独检查。

夹逼定理要求下界与上界具有相同的极限。利用 $s-1<\lfloor s\rfloor\le s$ 可得 $|t\lfloor1/t\rfloor-1|\le|t|$，故 $t\lfloor1/t\rfloor\to1$（$t\to0$）。若 $a_1,\ldots,a_m\ge0$、$M=\max_ka_k$，则
\[
    M\le\left(\sum_{k=1}^ma_k^p\right)^{1/p}\le m^{1/p}M
    \quad(p>0),
\]
所以 $p\to+\infty$ 时该表达式趋于 $M$。如果无法一次找到同极限的两端，也可对任意 $\eta>0$ 建立相差至多 $\eta$ 的渐近上、下界，最后再令 $\eta\downarrow0$。

\begin{example}
\label{ex:pei-convolution-means}
设 $x_n\to a$、$y_n\to b$，证明 $\displaystyle\frac1n\sum_{k=1}^nx_ky_{n+1-k}\to ab$。
\end{example}
\begin{solve}
令 $x_k=a+\alpha_k$、$y_k=b+\beta_k$，展开后除常数 $ab$ 外，出现两个趋零数列的算术平均，以及 $n^{-1}\sum_k\alpha_k\beta_{n+1-k}$。前两者由定理~\ref{thm:pei-weighted-cesaro} 趋于零。又因 $\alpha_k\to0$，存在 $M$ 使 $|\alpha_k|\le M$，故最后一项的绝对值不超过 $Mn^{-1}\sum_{k=1}^n|\beta_k|\to0$。这种替换把一般极限归结到零极限，并未交换随 $n$ 变化的求和与极限。
\end{solve}

由例~\ref{ex:pei-euler-constant}，有 $H_{2n}-H_n=\log2+o(1)$，从而
\begin{equation}
    \lim_{n\to\infty}\sum_{k=n+1}^{2n}\frac1k=\log2.
    \label{eq:pei-harmonic-block}
\end{equation}
另一种证明是逐项使用式~\eqref{eq:pei-log-bounds} 并求和，所得对数相消后同样夹逼到 $\log2$。同一个问题采用不同方法，关键在于识别它究竟具有平均结构、相消结构，还是积分和结构。

\subsection{微分、积分与级数方法的复习}

本小节引用后续微积分中的工具，适合在掌握对应定理后回看；这里只说明使用条件与关键计算，不把这些工具的完整证明作为本章的前提。\emph{L'Hospital 法则}要求 $f,g$ 在有关去心单侧邻域内可微，$g'$ 不为零，且 $f,g$ 同趋零，或者 $|g|\to+\infty$；若 $f'/g'\to L$（允许 $L=\pm\infty$），则 $f/g\to L$。分母趋于无穷大的版本不要求分子也趋于无穷大。乘积、差和幂式须先化成适用的商式，每次再次使用法则时都须重新核对条件。导数之比没有极限，并不能反推原商没有极限，例如 $(x+\sin x)/x\to1$，而导数之比 $1+\cos x$ 不收敛。

若 $f$ 在 $a$ 的邻域内有连续的 $m$ 阶导数，则带 Peano 余项的\emph{Taylor 公式}为
\[
    f(a+h)=\sum_{k=0}^m\frac{f^{(k)}(a)}{k!}h^k+o(h^m).
\]
当低阶项相消时，应展开到首个未抵消的阶次。例如
\[
    n\log\left(1+\frac1n\right)
    =1-\frac1{2n}+O(n^{-2}),\qquad
    n\left[e-\left(1+\frac1n\right)^n\right]\longrightarrow\frac e2.
\]
后一结论来自 $(1+1/n)^n=e\exp[-1/(2n)+O(n^{-2})]$，随后使用 $e^t-1\sim t$，所以余项在乘以 $n$ 后仍趋于零。

若 $f$ 在 $[a,b]$ 上 Riemann 可积，则等距分割的任意取样积分和趋于积分；特别，$[0,1]$ 上连续函数满足
\[
    \frac1n\sum_{k=1}^nf(k/n)\longrightarrow\int_0^1f(t)\,dt.
\]
取 $f(t)=1/(1+t)$，再得式~\eqref{eq:pei-harmonic-block}。对连乘先取对数，例如
\[
    \left[\prod_{k=1}^n\left(1+\frac kn\right)\right]^{1/n}
    \longrightarrow\exp\left(\int_0^1\log(1+t)\,dt\right)=\frac4e.
\]
\emph[caution]{若取样函数在端点无界，不能直接套用上述有界区间上的 Riemann 和结论}，必须另外控制端点贡献或改用其他方法。

设 $f\ge0$ 在 $[0,1]$ 上连续，$M=\max f$，则
\begin{equation}
    \lim_{p\to+\infty}\left(\int_0^1f(t)^p\,dt\right)^{1/p}=M.
    \label{eq:pei-lp-max}
\end{equation}
若 $M=0$，结论直接成立；若 $M>0$，对 $0<\eta<M$，连续性保证在某个正长度 $\ell_\eta$ 的子区间上 $f>M-\eta$，于是 $(M-\eta)\ell_\eta^{1/p}\le(\int f^p)^{1/p}\le M$。先令 $p\to\infty$，再令 $\eta\downarrow0$ 即得结论，端点处取到最大值也可使用单侧子区间。

收敛级数可用于控制通项或尾项。若 $\sum u_k$ 收敛，则 $u_k\to0$，且任意有限尾段 $\sum_{k=n+1}^{m(n)}u_k\to0$，其中 $m(n)\ge n$；后者由级数的 Cauchy 准则保证。若 $\sum_{n=2}^\infty|x_n-x_{n-1}|<+\infty$，则对 $m>n$ 有 $|x_m-x_n|\le\sum_{k=n+1}^m|x_k-x_{k-1}|\to0$，所以 $x_n$ 收敛。对于正数列，若 $u_{n+1}/u_n\to q<1$，则从某项起有 $u_{n+1}\le r u_n$，其中 $q<r<1$，几何估计已足以证明 $u_n\to0$，不必先证明整个级数收敛。

渐近增长比较也常用于选择估计。对固定 $\alpha>0$、$a>1$，当 $n\to\infty$ 时，
\[
    \log\log n\ll\log n\ll n^\alpha\ll a^n\ll n!\ll n^n,
\]
其中 $u_n\ll v_n$ 表示 $u_n/v_n\to0$。对数与幂、幂与指数的比较可由 L'Hospital 法则得到；$a^n/n!\to0$ 可用相邻比值 $a/(n+1)\to0$，而 $n!/n^n\le2^{-\lfloor n/2\rfloor}\to0$。各参数固定这一条件不可省略。

\subsection{周期平均与辅助函数}

\begin{example}
\label{ex:pei-periodic-average}
若 $f$ 是周期为 $T>0$ 的连续函数，证明
\[
    \lim_{x\to+\infty}\frac1x\int_0^xf(t)\,dt
    =\frac1T\int_0^Tf(t)\,dt.
\]
\end{example}
\begin{solve}
写成 $x=nT+r$，其中 $0\le r<T$。连续周期函数有界，设 $|f|\le M$，则
\[
    \int_0^xf=n\int_0^Tf+\int_0^rf,
    \qquad\left|\int_0^rf\right|\le MT.
\]
除以 $x$，由 $n/x\to1/T$ 及有界余项除以 $x$ 趋零，得到结论。这里长时间平均由完整周期决定，末尾不足一周期的贡献有统一上界。
\end{solve}

\begin{example}
\label{ex:pei-integrating-factor}
设 $f\in C^1([0,+\infty))$，且 $f(x)+f'(x)\to0$，证明 $f(x)\to0$。
\end{example}
\begin{solve}
令 $h=f+f'$。由 $(e^xf(x))'=e^xh(x)$，对任意固定 $A\ge0$ 和 $x\ge A$，有
\[
    f(x)=e^{A-x}f(A)+\int_A^xe^{t-x}h(t)\,dt.
\]
给定 $\varepsilon>0$，先选 $A$ 使 $t\ge A$ 时 $|h(t)|<\varepsilon/2$，积分的绝对值便小于 $\varepsilon/2$；再固定此 $A$，令 $x$ 足够大，使首项绝对值小于 $\varepsilon/2$，即得结论。
\end{solve}

类似地，若 $f\in C^1(\mathbb{R})$ 有界且 $|f-f'|\le1$，令 $h=f-f'$，对 $R>x$ 积分 $ (e^{-t}f(t))'=-e^{-t}h(t)$，得到
\[
    f(x)=e^{x-R}f(R)+\int_x^Re^{x-t}h(t)\,dt.
\]
令 $R\to+\infty$，有界性消去首项，故 $|f(x)|\le\int_x^\infty e^{x-t}\,dt=1$。辅助函数的选择来自希望把 $f$ 与 $f'$ 合并成一个乘积的导数，而所需无穷积分由指数核的可积性保证。

% 第 1.4 节：PDF 第 46--55 页。
\section{Stolz 公式}
\label{sec:pei-stolz}

\subsection{两种数列形式及其证明}

\begin{theorem}[Stolz 公式]
\label{thm:pei-stolz}
设 $\{a_n\}$、$\{b_n\}$ 为实数列。如果 $b_n$ 严格递增且 $b_n\to+\infty$，并且
\begin{equation}
    \frac{a_{n+1}-a_n}{b_{n+1}-b_n}\longrightarrow L,
    \label{eq:pei-stolz-differences}
\end{equation}
其中 $L\in\mathbb{R}\cup\{\pm\infty\}$，则 $a_n/b_n\to L$。另一个版本是：若 $b_n>0$ 严格递减到零，$a_n\to0$，且式~\eqref{eq:pei-stolz-differences} 成立，则仍有 $a_n/b_n\to L$。
\end{theorem}
\begin{proof}
先设第一种情形下 $L$ 有限。给定 $\varepsilon>0$，存在 $N$ 使 $k\ge N$ 时
\[
    (L-\varepsilon)(b_{k+1}-b_k)
    \le a_{k+1}-a_k
    \le(L+\varepsilon)(b_{k+1}-b_k).
\]
从 $k=N$ 累加到 $n-1$，得 $|(a_n-a_N)-L(b_n-b_N)|\le\varepsilon(b_n-b_N)$。除以充分后期为正的 $b_n$，再令 $n\to\infty$，有限首段的贡献趋零，故 $\limsup|a_n/b_n-L|\le\varepsilon$；因 $\varepsilon$ 任意，结论成立。若 $L=+\infty$，对任意 $M>0$ 从某项起累加 $a_{k+1}-a_k\ge M(b_{k+1}-b_k)$，得到 $\liminf a_n/b_n\ge M$；$L=-\infty$ 时改看 $-a_n$。

第二种情形把差写成 $a_k-a_{k+1}$、$b_k-b_{k+1}>0$。有限 $L$ 时从 $k=n$ 累加到 $m-1$，再利用 $a_m,b_m\to0$，得到 $|a_n-Lb_n|\le\varepsilon b_n$，除以 $b_n$ 即可。无穷极限时作相同的单侧估计。
\end{proof}

第一种形式不要求分子趋于无穷大；第二种形式则必须同时要求分子、分母趋零。公式提供充分条件，不提供逆命题：例如 $a_n=(-1)^n$、$b_n=n$ 时 $a_n/b_n\to0$，但相邻差之比不收敛。使用时应先确认分母的单调性及去向，再计算相邻差，不能只凭式子看起来像不定式就套用。

\subsection{求和、连乘与比值的应用}

对固定实数 $p>-1$，令 $a_n=\sum_{k=1}^nk^p$、$b_n=n^{p+1}$，由中值定理有 $(n+1)^{p+1}-n^{p+1}\sim(p+1)n^p$，故
\begin{equation}
    \frac{1^p+2^p+\cdots+n^p}{n^{p+1}}\longrightarrow\frac1{p+1}.
    \label{eq:pei-power-sums}
\end{equation}
加权平均定理也可由 Stolz 公式验证：若 $p_n>0$，则 $P_n$ 严格递增，且相邻分子、分母差的比正好是 $a_{n+1}$；若部分 $p_n=0$，则直接使用定理~\ref{thm:pei-weighted-cesaro}，或去掉零权重后再应用公式。

\begin{example}
\label{ex:pei-factorial-root}
证明 $\displaystyle\frac{(n!)^{1/n}}n\to e^{-1}$。
\end{example}
\begin{solve}
对原式取对数，得到 $A_n/n$，其中 $A_n=\log(n!)-n\log n$。其相邻差为
\[
    A_{n+1}-A_n=-n\log\left(1+\frac1n\right)\longrightarrow-1.
\]
以 $b_n=n$ 应用 Stolz 公式，得到 $A_n/n\to-1$，再取指数即得结论。此法避免了把端点处无界的 $\log t$ 直接当作普通 Riemann 和函数使用，也不需要预先引用 Stirling 公式。
\end{solve}

若 $u_n>0$ 且 $u_{n+1}/u_n\to q>0$，对 $\log u_n/n$ 使用 Stolz 公式，得到 $u_n^{1/n}\to q$。若相邻比值趋于零，则对任意 $r>0$，充分后期有 $u_{n+1}\le r u_n$，故 $\limsup u_n^{1/n}\le r$，从而根式极限为零；相邻比值趋于 $+\infty$ 时类似得到根式趋于 $+\infty$。第~\ref{sec:pei-sequence-limsup} 节将把这些结论推广到比值不收敛的情形。

\subsection{函数极限的差分版本}

\begin{proposition}
\label{prop:pei-function-stolz}
设 $T>0$，$f$ 在 $[a,+\infty)$ 的每个有界闭子区间上有界，$g$ 严格递增且 $g(x)\to+\infty$。若
\[
    \frac{f(x+T)-f(x)}{g(x+T)-g(x)}\longrightarrow L
    \qquad(x\to+\infty),
\]
则 $f(x)/g(x)\to L$，其中 $L$ 可为有限实数或 $\pm\infty$。另若 $g(x)>0$、$g(x+T)<g(x)$，$f(x),g(x)\to0$，且同一差分比趋于 $L$，则也有 $f(x)/g(x)\to L$。
\end{proposition}
\begin{proof}
第一种情形下，选取充分大的 $A$，把每个足够大的 $x$ 唯一写成 $x=t+nT$，其中 $t\in[A,A+T)$。沿 $t,t+T,\ldots,t+nT$ 累加差分不等式，便得数列版本中的估计。关键是它对 $t$ 一致成立：$f(t)$ 在该初始区间上有界，单调的 $g(t)$ 也有界，而 $g(x)\to+\infty$，故所有首段误差除以 $g(x)$ 后一致趋零。第二种情形沿每条等距点列把估计累加至无穷，利用 $f,g\to0$ 消去尾端，即得与数列版本相同的界。
\end{proof}

特别，若 $f$ 局部有界，且 $f(x+1)-f(x)\to A$，则取 $g(x)=x$ 得 $f(x)/x\to A$。局部有界保证不同小数部分对应的首段值受到统一控制，不能仅对整数点列证明后就把结论扩展到全体实变量。

% 第 1.5 节：PDF 第 55--78 页。
\section{递推形式的极限}
\label{sec:pei-recursive-limits}

\subsection{不动点、单调性与收敛依据}

递推式 $x_{n+1}=f(x_n)$ 中满足 $f(c)=c$ 的点称为\emph{不动点}（fixed point）。若已知 $x_n\to A$ 且 $f$ 在 $A$ 处连续，才能由递推式推出 $A=f(A)$；解出这个方程只得到候选值。实际解题时可以先求候选点，以便建立界或估计误差，但随后必须独立验证收敛。还要注意，方程即使只有一个不动点，迭代也可能不收敛，例如 $f(x)=1-x$ 在 $[0,1]$ 上的不动点唯一，而从 $x_0=0$ 出发得到 $0,1,0,1,\ldots$。

\begin{example}
\label{ex:pei-square-root-recursion}
设 $x_0=1$、$x_{n+1}=\sqrt{2x_n}$，求其极限。
\end{example}
\begin{solve}
由归纳法，$1\le x_n<2$。在这个范围内 $2x_n\ge x_n^2$，故 $x_{n+1}\ge x_n$。由单调有界原理，$x_n\to A\in[1,2]$；再由连续性得到 $A=\sqrt{2A}$，候选根为 $0,2$，其中只有 $2$ 位于该范围，所以极限为 $2$。也可直接写出 $x_n=2^{1-2^{-n}}$，独立核查结果。
\end{solve}

对单调不减的连续函数，若已找到不动点 $c$，且 $x<c$ 时 $x\le f(x)\le c$、$x>c$ 时 $c\le f(x)\le x$，则从相应一侧出发的迭代单调有界，极限是该侧范围内的某个不动点；若该范围内不动点唯一，则极限必为 $c$。例如，已知存在 $c$ 且 $0<f'(x)<1$ 时，$f$ 递增、$f(x)-x$ 严格递减，因而满足上述指向 $c$ 的关系。这里不需要把逐点的 $f'(x)<1$ 错当成存在统一常数 $q<1$ 的导数上界。

单调性也允许可求和的扰动。若 $x_n\ge m$，且 $x_{n+1}\le x_n+\varepsilon_n$，其中 $\varepsilon_n\ge0$、$\sum_n\varepsilon_n<+\infty$，令 $r_n=\sum_{k=n}^\infty\varepsilon_k$，则 $x_n+r_n$ 单调不增且有下界，所以收敛；又 $r_n\to0$，故 $x_n$ 收敛。它把“每一步略有增加”的问题转成真正的单调数列，条件是全部向上扰动的总量有限。

\subsection{压缩估计与不变区间}

\begin{theorem}[相邻差的压缩判别]
\label{thm:pei-contractive-differences}
若存在固定 $q\in[0,1)$，使所有 $n\ge1$ 满足 $|x_{n+1}-x_n|\le q|x_n-x_{n-1}|$，则 $x_n$ 收敛。若其极限为 $A$，则
\begin{equation}
    |A-x_n|\le\frac{q^n}{1-q}|x_1-x_0|\qquad(n\ge1).
    \label{eq:pei-contraction-error}
\end{equation}
\end{theorem}
\begin{proof}
归纳得 $|x_{k+1}-x_k|\le q^k|x_1-x_0|$。对 $m>n$ 求和，
\[
    |x_m-x_n|\le\sum_{k=n}^{m-1}|x_{k+1}-x_k|
    \le\frac{q^n}{1-q}|x_1-x_0|.
\]
右端随 $n$ 趋零且与 $m$ 无关，故由 Cauchy 准则得收敛；再令 $m\to\infty$ 即得误差界。$q=0$ 时从 $x_1$ 起数列恒定，结论同样成立。
\end{proof}

\begin{theorem}[闭区间上的压缩迭代]
\label{thm:pei-contraction-map}
设 $I$ 为非空闭区间，$f:I\to I$，且存在 $q\in[0,1)$ 使
\[
    |f(x)-f(y)|\le q|x-y|\qquad(x,y\in I).
\]
则 $f$ 在 $I$ 上有唯一不动点 $c$，任取 $x_0\in I$ 产生的迭代 $x_{n+1}=f(x_n)$ 都趋于 $c$，且 $|x_n-c|\le q^n|x_0-c|$。
\end{theorem}
\begin{proof}
$f(I)\subseteq I$ 保证所有迭代有定义并留在 $I$ 中，相邻差满足定理~\ref{thm:pei-contractive-differences} 的条件，故收敛。闭性保证极限 $c\in I$，压缩条件保证 $f$ 连续，因而 $f(c)=c$。若 $d$ 也是不动点，则 $|c-d|\le q|c-d|$，只能有 $c=d$。把 $y=c$ 代入压缩条件并归纳，即得最后的误差界。
\end{proof}

若 $f$ 在区间内可微，且 $|f'|\le q<1$，中值定理便给出压缩条件，但仍须先验证 $f(I)\subseteq I$。例如取 $I=[a-r,a+r]$，若 $|f(a)-a|\le(1-q)r$ 且 $|f'|\le q$，则 $|f(x)-a|\le q|x-a|+(1-q)r\le r$，不变性成立。仅有 $|f(x)-f(y)|\le|x-y|$，即\emph{非扩张性}（nonexpansiveness），不能使用上述严格压缩结论。

\begin{example}
\label{ex:pei-averaged-nonexpansive}
设 $f:[a,b]\to[a,b]$ 满足 $|f(x)-f(y)|\le|x-y|$。证明平均迭代 $x_{n+1}=[x_n+f(x_n)]/2$ 收敛到 $f$ 的一个不动点。
\end{example}
\begin{solve}
令 $G(x)=[x+f(x)]/2$，则 $G$ 连续并把 $[a,b]$ 映入自身。若 $x<y$，有 $G(y)-G(x)=[y-x+f(y)-f(x)]/2\ge0$，所以 $G$ 单调不减。由 $x_1\ge x_0$ 或 $x_1\le x_0$ 分别归纳出整个迭代单调，且始终有界，故收敛到某个 $c$。由 $c=G(c)$ 得 $f(c)=c$。本题使用的是一维次序和平均后的单调性，不需要严格压缩，也不保证不动点唯一。
\end{solve}

\subsection{交替迭代、通项与变量替换}

当 $f$ 连续、单调递减且把有界闭区间映入自身时，$F=f\circ f$ 单调不减。于是偶数子列和奇数子列分别是 $F$ 的迭代，均单调有界，极限 $\alpha,\beta$ 满足 $\alpha=f(\beta)$、$\beta=f(\alpha)$。只有进一步证明 $\alpha=\beta$，才能断定原列收敛；否则可能形成二周期。几何上，交替把 $x_n$ 送到 $f(x_n)$，再把函数值视为下一次自变量，可以帮助猜测两子列的方向，但图形观察不能替代这些不等式与极限论证。

\begin{example}
\label{ex:pei-golden-recursion}
设 $x_0=1$，$x_{n+1}=1/(1+x_n)$，求极限。
\end{example}
\begin{solve}
所有项都在 $[1/2,1]$ 中，且 $f(x)=1/(1+x)$ 在此区间上满足 $|f'(x)|\le4/9<1$，故定理~\ref{thm:pei-contraction-map} 保证收敛。极限满足 $A^2+A-1=0$，由正性得 $A=(\sqrt5-1)/2$。若用奇、偶子列法，则先得到 $\alpha=1/(1+\beta)$、$\beta=1/(1+\alpha)$，两式相减给出 $\alpha=\beta$，随后求出相同的根。若只在较大的 $[0,1]$ 上估计导数，所得上界为一；应缩小到实际不变区间，或改用子列论证。
\end{solve}

对于线性递推，写出通项或找不变量通常更合适。例如 $x_0=a$、$x_1=b$、$x_{n+2}=(x_{n+1}+x_n)/2$ 满足 $x_{n+2}-x_{n+1}=-(x_{n+1}-x_n)/2$，所以
\[
    x_n=\frac{a+2b}{3}+\frac{2(a-b)}3\left(-\frac12\right)^n
    \longrightarrow\frac{a+2b}{3}.
\]
在递推式中直接代入极限只得到 $A=A$，无法识别这个依赖初值的数值；差分通项保留了所需信息。

\begin{proposition}[二阶线性递推的比值]
\label{prop:pei-second-order-ratio}
设 $p,q>0$，$u_{n+2}=pu_{n+1}+qu_n$，$u_0,u_1\in\mathbb{R}$。记
\[
    \alpha=\frac{p+\sqrt{p^2+4q}}2,\qquad
    \beta=\frac{p-\sqrt{p^2+4q}}2.
\]
若 $u_1\ne\beta u_0$，则从某项起 $u_n\ne0$，且 $u_{n+1}/u_n\to\alpha$。若 $u_1=\beta u_0$、$u_0\ne0$，则该比值恒为 $\beta$；若 $u_0=u_1=0$，则比值无定义。
\end{proposition}
\begin{proof}
两个特征根不同，直接代入并核对初值可得 $u_n=C\alpha^n+D\beta^n$，其中 $C=(u_1-\beta u_0)/(\alpha-\beta)$、$D=u_0-C$。因 $\alpha>0>\beta$ 且 $|\beta|<\alpha$，当 $C\ne0$ 时除以 $C\alpha^n$，余项 $(D/C)(\beta/\alpha)^n\to0$，所以分母最终非零且比值趋于 $\alpha$。$C=0$ 时则只剩 $D\beta^n$。
\end{proof}

变量替换可以把陌生递推化为已有类型。例如 $x_{n+1}=2+1/x_n$、$x_0\ge2$ 时，$[2,5/2]$ 从第一次迭代后是不变区间，导数绝对值不超过 $1/4$，故极限为 $1+\sqrt2$。一般的分式递推也可先寻找不动点，再考察 $x_n-c$ 或 $1/x_n$ 的递推；替换时须核对分母始终不为零。

\begin{example}
\label{ex:pei-conjugate-fractional-part}
以 $\{x\}=x-\lfloor x\rfloor$ 表示小数部分，求 $\displaystyle\lim_{n\to\infty}\{(2+\sqrt3)^n\}$。
\end{example}
\begin{solve}
由二项式展开，$s_n=(2+\sqrt3)^n+(2-\sqrt3)^n$ 是整数，因为奇数次的根式项相消，偶数次项为整数的两倍。又 $r_n=(2-\sqrt3)^n\in(0,1)$ 且趋零，所以 $(2+\sqrt3)^n=s_n-r_n$，其小数部分为 $1-r_n\to1$。小数部分始终小于一，不妨碍其极限等于一；这里利用共轭表达式识别了相邻整数。
\end{solve}

\subsection{求根迭代与耦合递推}

\begin{example}
\label{ex:pei-newton-pth-root}
设 $a>0$、整数 $p\ge2$、$x_0\ge a^{1/p}$，定义
\[
    x_{n+1}=\frac1p\left((p-1)x_n+\frac{a}{x_n^{p-1}}\right).
\]
证明 $x_n\to a^{1/p}$。
\end{example}
\begin{solve}
对 $p-1$ 个 $x_n$ 与一个 $a/x_n^{p-1}$ 应用算术平均与几何平均不等式，得 $x_{n+1}\ge a^{1/p}$。若 $x_n\ge a^{1/p}$，则
\[
    x_{n+1}-x_n=\frac{a-x_n^p}{p x_n^{p-1}}\le0.
\]
故归纳得数列单调不增且有正下界。令极限为 $A$，代入递推式可得 $A^p=a$，由 $A>0$ 得所求结论。正下界同时保证取极限时分母不趋零。
\end{solve}

更一般的\emph{Newton 迭代}（Newton iteration）为 $x_{n+1}=x_n-f(x_n)/f'(x_n)$。一个可核查的充分条件是：$f\in C^2([a,b])$，$f'>0$、$f''\ge0$，$f(a)\le0\le f(b)$，并从 $x_0=b$ 出发。此时连续性和严格单调性保证唯一零点 $c\in[a,b]$。对 $x\ge c$，由导数不减有 $0\le f(x)=\int_c^xf'(t)\,dt\le f'(x)(x-c)$，所以 $c\le x-f(x)/f'(x)\le x$。迭代因而单调下降且留在 $[c,b]$；若极限为 $A$，连续性及 $f'(A)>0$ 给出 $f(A)=0$，所以 $A=c$。离开这些条件时，切线图示本身不构成收敛保证。

\begin{example}
\label{ex:pei-agm-recursion}
设 $a_0\ge b_0>0$，定义 $a_{n+1}=(a_n+b_n)/2$、$b_{n+1}=\sqrt{a_nb_n}$。证明两列收敛到同一个数。
\end{example}
\begin{solve}
由算术平均与几何平均不等式，
\[
    b_n\le b_{n+1}\le a_{n+1}\le a_n.
\]
故两列单调有界，分别趋于 $A,B$，且 $A\ge B\ge b_0>0$。在算术平均递推中取极限，$A=(A+B)/2$，于是 $A=B$。这个公共值称为两初值的\emph{算术几何平均}（arithmetic-geometric mean），一般不能仅凭递推关系把它写成初值的一个简单代数式。
\end{solve}

二元递推也可用同样的次序思想。设连续正函数 $F(x,y)$ 分别对两个变量单调不减，并存在 $c>0$，使 $F(c,c)=c$，$0<t<c$ 时 $t<F(t,t)$，$t>c$ 时 $F(t,t)<t$。对 $x_0=x_1=d>0$、$x_{n+2}=F(x_{n+1},x_n)$，若 $d<c$，则由归纳法得到数列单调不减且不超过 $c$；若 $d>c$，则单调不增且不低于 $c$。极限 $A>0$ 满足 $A=F(A,A)$，上述严格符号条件又排除了 $A\ne c$，故两种情形都趋于 $c$。分别单调性保证对角线以外的递推也服从这些界，不能只检查 $F(t,t)$ 的图像。

\subsection{直接误差法与收敛速度}

如果已猜出候选极限 $c$，也可直接证明误差 $e_n=x_n-c$ 趋零，无须先证明 $x_n$ 单调。例如若 $|e_{n+1}|\le q|e_n|$、$0\le q<1$，则 $|e_n|\le q^n|e_0|$。对于二阶递推，若 $|e_{n+1}|\le q\max\{|e_n|,|e_{n-1}|\}$，令 $M_k=\max\{|e_{2k}|,|e_{2k+1}|\}$，连续使用两次估计便有 $M_{k+1}\le qM_k$，仍可得到收敛。关键是找到一个固定的严格小于一的因子，而不是随指标接近一的数。

\begin{proposition}[趋零递推的渐近速度]
\label{prop:pei-recursion-rate}
设 $x_n>0$、$x_n\to0$，存在常数 $c,\alpha>0$ 使
\[
    x_{n+1}=x_n-cx_n^{1+\alpha}+o(x_n^{1+\alpha}).
\]
则 $n x_n^\alpha\to1/(\alpha c)$。
\end{proposition}
\begin{proof}
令 $y_n=x_n^{-\alpha}$，则
\[
    y_{n+1}-y_n
    =x_n^{-\alpha}\left[
       (1-cx_n^\alpha+o(x_n^\alpha))^{-\alpha}-1\right]
    \longrightarrow\alpha c.
\]
这里使用零点处的一阶展开；对 $y_n/n$ 应用 Stolz 公式，得 $y_n/n\to\alpha c$，取倒数即得结论。
\end{proof}

\begin{example}
\label{ex:pei-sine-recursion-rate}
设 $0<x_0<\pi/2$、$x_{n+1}=\sin x_n$，证明 $x_n\to0$，并求其渐近速度。
\end{example}
\begin{solve}
由于 $0<\sin x<x$（$0<x<\pi/2$），数列递减且为正，所以有极限 $A\ge0$。由 $A=\sin A$，在 $[0,x_0]$ 上只能有 $A=0$。再使用 $\sin x=x-x^3/6+o(x^3)$，在命题~\ref{prop:pei-recursion-rate} 中取 $\alpha=2$、$c=1/6$，得到 $nx_n^2\to3$，即 $\sqrt{n/3}\,x_n\to1$。本题的不动点处导数为一，严格压缩估计不能直接给出这一幂次速度。
\end{solve}

同理，若 $0<a_0<1$、$a_{n+1}=a_n(1-a_n)$，先由单调有界及不动点方程证明 $a_n\to0$，再取 $\alpha=c=1$，得到 $na_n\to1$。收敛性和速度分成两个步骤：前者排除非零极限，后者利用首个非零误差项选择合适的倒数幂，再交给 Stolz 公式。

% 第 1.6 节：PDF 第 78--89 页。
\section{序列的上、下极限}
\label{sec:pei-sequence-limsup}

\subsection{尾部确界与量词刻画}

对不收敛的数列，仅说“极限不存在”会丢掉振荡范围的信息。为统一处理无界情形，以下允许上、下极限取扩展实数 $\overline{\mathbb{R}}=\mathbb{R}\cup\{-\infty,+\infty\}$，但不把 $+\infty-\infty$ 或 $0\cdot\infty$ 视为已有定义的运算。

\begin{definition}
\label{def:pei-sequence-limsup}
对实数列 $\{x_n\}$，令 $M_n=\sup_{k\ge n}x_k$、$m_n=\inf_{k\ge n}x_k$，允许无界集合的确界为 $\pm\infty$。定义其\emph{上极限}（limit superior）和\emph{下极限}（limit inferior）为
\begin{equation}
    \limsup_{n\to\infty}x_n=\lim_{n\to\infty}M_n=\inf_nM_n,
    \qquad
    \liminf_{n\to\infty}x_n=\lim_{n\to\infty}m_n=\sup_nm_n.
    \label{eq:pei-tail-extrema}
\end{equation}
因 $M_n$ 单调不增、$m_n$ 单调不减，上述扩展实数极限总存在。
\end{definition}

记 $\ell=\liminf x_n$、$u=\limsup x_n$，总有 $\ell\le u$。若 $u$ 有限，它等价地由两个条件刻画：任意 $\varepsilon>0$ 时最终有 $x_n<u+\varepsilon$；任意 $\varepsilon>0$、任意 $N$ 之后仍能找到 $n>N$ 使 $x_n>u-\varepsilon$。第一条限制最终上界，第二条保证数列不断逼近这个边界。下极限的对应条件是最终有 $x_n>\ell-\varepsilon$，且任意后段都存在 $x_n<\ell+\varepsilon$。这些条件中的严格不等号与非严格不等号可通过缩小 $\varepsilon$ 相互转换。

例如 $x_n=(-1)^n+1/n$ 的上极限为一，下极限为负一。若 $x_{2n}=0$、$x_{2n-1}=n$，则上极限为 $+\infty$、下极限为零；上极限为 $+\infty$ 只说明没有有限上界，不等于整个数列趋于 $+\infty$。后者要求任意 $M$ 之上包含全部充分后期的项，也即 $\liminf x_n=+\infty$。

\subsection{子列解释与普通极限}

\begin{theorem}
\label{thm:pei-limsup-subsequences}
实数列的上极限是其全部广义子列极限的最大者，下极限是最小者；“广义”允许极限为 $\pm\infty$。特别，存在子列分别趋于上极限与下极限，任意子列 $x_{n_k}$ 都满足
\[
    \liminf_nx_n\le\liminf_kx_{n_k}
       \le\limsup_kx_{n_k}\le\limsup_nx_n.
\]
\end{theorem}
\begin{proof}
先设上极限 $u$ 有限。由尾部逼近性质，可递归选取 $n_j>n_{j-1}$ 且 $x_{n_j}>u-1/j$；最终上界又保证 $x_{n_j}\to u$。任何收敛子列的极限均不超过 $u$，否则它的充分后期会超过原列的某个最终上界。若 $u=+\infty$，逐次选取 $n_j>n_{j-1}$ 使 $x_{n_j}>j$；若 $u=-\infty$，则全列已趋于 $-\infty$。下极限对 $-x_n$ 使用同样论证。子列尾部包含在原列的相应尾部中，故其确界范围不会超出原列的最终范围，得到所列不等式。
\end{proof}

\begin{corollary}
\label{cor:pei-limsup-convergence}
对 $A\in\overline{\mathbb{R}}$，有 $x_n\to A$ 当且仅当 $\liminf x_n=\limsup x_n=A$；普通收敛还要求 $A$ 有限。
\end{corollary}
\begin{proof}
有限 $A$ 时，最终上、下界把 $x_n$ 夹在 $(A-\varepsilon,A+\varepsilon)$ 中；$A=\pm\infty$ 时对应单侧最终界即为趋于无穷大的定义。反向直接由极限定义和尾部确界得到。
\end{proof}

\subsection{运算不等式及符号条件}

\begin{proposition}
\label{prop:pei-limsup-sum}
设 $x_n,y_n$ 有界，记 $\ell_x=\liminf x_n$、$u_x=\limsup x_n$，$\ell_y,u_y$ 类似定义，则
\begin{equation}
    \begin{aligned}
        \ell_x+\ell_y&\le\liminf(x_n+y_n)
          \le\min\{\ell_x+u_y,u_x+\ell_y\},\\
        \max\{\ell_x+u_y,u_x+\ell_y\}
          &\le\limsup(x_n+y_n)\le u_x+u_y.
    \end{aligned}
    \label{eq:pei-limsup-sum-bounds}
\end{equation}
\end{proposition}
\begin{proof}
任意 $\varepsilon>0$ 时，充分后期有 $\ell_x-\varepsilon<x_n<u_x+\varepsilon$，$\ell_y-\varepsilon<y_n<u_y+\varepsilon$；相加并令 $\varepsilon\downarrow0$，得到最外侧的界。取 $x_{n_k}\to\ell_x$ 的子列，利用 $y_{n_k}$ 的最终上界，可得 $\liminf(x_n+y_n)\le\ell_x+u_y$；交换两列得到另一个上界。对两列取负并用 $\limsup(-z_n)=-\liminf z_n$，得到第二行的内侧界。
\end{proof}

一般不能把这些界改成等式。例如 $x_n=(-1)^n$、$y_n=-(-1)^n$ 时，和恒为零，而两个上极限之和为二。若其中一列 $x_n\to a\in\mathbb{R}$，则上下夹界合并，得到
\[
    \limsup(x_n+y_n)=a+\limsup y_n,\qquad
    \liminf(x_n+y_n)=a+\liminf y_n.
\]
此处也允许 $y_n$ 无界，因为加上趋于零的误差不会改变扩展实数的最终界。

若 $x_n,y_n$ 均非负且有界，同样的确界估计给出
\[
    \ell_x\ell_y\le\liminf(x_ny_n)
      \le\min\{\ell_xu_y,u_x\ell_y\},\qquad
    \max\{\ell_xu_y,u_x\ell_y\}
      \le\limsup(x_ny_n)\le u_xu_y.
\]
证明内侧界时选取达到相应上、下极限的子列，非负性保证乘法保持不等号方向。若允许负值，则必须分别处理符号；特别，当 $x_n\to a$ 且 $y_n$ 有界时，$a>0$ 保持上、下极限，$a<0$ 交换上、下极限，$a=0$ 时乘积趋零。若 $\liminf y_n>0$ 且 $y_n$ 有界，则 $\limsup(1/y_n)=1/\liminf y_n$、$\liminf(1/y_n)=1/\limsup y_n$，可据此处理商。无界情形应另行检查是否出现未定义的扩展实数运算。

\subsection{相邻比值、根式与次可加性}

\begin{proposition}
\label{prop:pei-ratio-root-limsup}
对任意正数列 $a_n$，有
\begin{equation}
    \liminf\frac{a_{n+1}}{a_n}
    \le\liminf a_n^{1/n}
    \le\limsup a_n^{1/n}
    \le\limsup\frac{a_{n+1}}{a_n}.
    \label{eq:pei-ratio-root-limsup}
\end{equation}
各量允许取零或 $+\infty$。
\end{proposition}
\begin{proof}
记右端为 $U$。若 $U<+\infty$，任取 $r>U$，充分后期有 $a_{n+1}\le r a_n$，故 $a_n\le a_Nr^{n-N}$。开 $n$ 次方并取上极限，得 $\limsup a_n^{1/n}\le r$，再令 $r\downarrow U$；若 $U=+\infty$，该侧不等式直接成立。下界类似：若 $L=\liminf a_{n+1}/a_n>0$，对每个 $0<r<L$ 作反向的几何估计，得到 $\liminf a_n^{1/n}\ge r$，再令 $r\uparrow L$；$L=0$ 时由正性得到结论，$L=+\infty$ 时可取任意 $r>0$。
\end{proof}

\begin{theorem}[次可加数列的平均极限]
\label{thm:pei-subadditive}
设实数列 $a_n$ 满足\emph{次可加性}（subadditivity），即 $a_{m+n}\le a_m+a_n$（$m,n\ge1$）。则
\begin{equation}
    \lim_{n\to\infty}\frac{a_n}{n}
       =\inf_{m\ge1}\frac{a_m}{m},
    \label{eq:pei-subadditive}
\end{equation}
右端允许为 $-\infty$。
\end{theorem}
\begin{proof}
固定 $m\ge1$，把 $n=qm+r$，$0\le r<m$，并约定 $a_0=0$。反复使用次可加性，得 $a_n\le qa_m+a_r$，其中 $r=0$ 时右端就是 $qa_m$。因有限个余数项有统一上界，除以 $n$ 后取上极限，得到 $\limsup a_n/n\le a_m/m$。若 $I=\inf_ma_m/m$ 有限，再对 $m$ 取下确界，并结合每项 $a_n/n\ge I$，即得上下极限相等；若 $I=-\infty$，上述上极限估计可小于任意实数，说明 $a_n/n\to-\infty$。
\end{proof}

这个证明的核心是固定块长后控制有限种余数，再让总长度趋于无穷。它也展示了上、下极限的一般用法：不预设收敛，先分别估计，再证明两者必须重合。

% 第 1.7 节：PDF 第 89--95 页。
\section{函数的上、下极限}
\label{sec:pei-function-limsup}

\subsection{子极限与去心邻域确界}

设 $f:E\to\mathbb{R}$，$a$ 是 $E$ 的聚点。若存在 $t_n\in E\setminus\{a\}$，$t_n\to a$ 且 $f(t_n)\to L\in\overline{\mathbb{R}}$，则称 $L$ 为 $f$ 在 $a$ 处的\emph{子极限}（cluster value）。它记录沿某条趋近途径得到的极限；上、下极限则分别记录所有这些途径中的最大、最小边界。

\begin{definition}
\label{def:pei-function-limsup}
令 $E_\delta=\{x\in E:0<|x-a|<\delta\}$，它对每个 $\delta>0$ 都非空。定义
\begin{equation}
    \begin{aligned}
        \limsup_{\substack{x\to a\\x\in E}}f(x)
          &=\lim_{\delta\downarrow0}\sup_{x\in E_\delta}f(x)
           =\inf_{\delta>0}\sup_{x\in E_\delta}f(x),\\
        \liminf_{\substack{x\to a\\x\in E}}f(x)
          &=\lim_{\delta\downarrow0}\inf_{x\in E_\delta}f(x)
           =\sup_{\delta>0}\inf_{x\in E_\delta}f(x).
    \end{aligned}
    \label{eq:pei-function-tail-extrema}
\end{equation}
确界和极限允许取扩展实数值。邻域缩小时上确界不增、下确界不减，所以两者总存在。
\end{definition}

有限上极限 $U$ 的量词刻画是：对任意 $\varepsilon>0$，存在 $\delta>0$，使全部 $x\in E_\delta$ 满足 $f(x)<U+\varepsilon$；同时对任意 $\varepsilon>0$ 和任意 $\delta>0$，总有某个 $x\in E_\delta$ 满足 $f(x)>U-\varepsilon$。下极限把两种不等号反向。由第二条递归选取 $|t_n-a|<1/n$ 的点，并结合第一条，可使 $f(t_n)\to U$；其余子极限不超过 $U$。对无穷边界以 $n$、$-n$ 替代 $U\pm1/n$ 作同样选取，因而上、下极限确实分别是最大、最小广义子极限。

设 $t_n\to a$、$t_n\in E\setminus\{a\}$，则
\[
    \liminf_{x\to a}f(x)\le\liminf_nf(t_n)
       \le\limsup_nf(t_n)\le\limsup_{x\to a}f(x).
\]
这给出函数与数列两种上、下极限之间的联系，且两端均可由适当选取趋近数列达到。普通函数极限存在的充要条件是这两个边界重合且有限；若重合于 $\pm\infty$，则对应趋于该无穷大。

\subsection{单侧形式与振荡例子}

若 $a$ 同时是 $E$ 左侧和右侧的聚点，在 $E_\delta$ 中分别增加 $x<a$ 或 $x>a$ 的限制，便得到左、右上极限和下极限。去心邻域是两侧邻域的并，因此
\[
    \begin{aligned}
        \limsup_{x\to a}f(x)
          &=\max\left\{\limsup_{x\to a-}f(x),
                       \limsup_{x\to a+}f(x)\right\},\\
        \liminf_{x\to a}f(x)
          &=\min\left\{\liminf_{x\to a-}f(x),
                       \liminf_{x\to a+}f(x)\right\}.
    \end{aligned}
\]
若定义域只从一侧积聚，应只讨论该侧，不能给另一侧虚设一个极限。$x\to+\infty$ 时，把去心邻域换成尾部 $E\cap(R,+\infty)$ 即得对应定义；$x\to-\infty$ 类似。

\begin{example}
\label{ex:pei-oscillation-cluster-values}
求 $f(x)=\sin(1/x)$ 在 $x\to0$ 时，以及 $D(x)=\mathbf{1}_{\mathbb{Q}}(x)$ 在任意点处的上、下极限和子极限。
\end{example}
\begin{solve}
前者始终在 $[-1,1]$ 中。任取 $c\in[-1,1]$，选 $\theta$ 使 $\sin\theta=c$，再令 $t_n=1/(2\pi n+\theta)$（去掉有限个分母为零的指标），则 $t_n\to0$ 且 $f(t_n)=c$。所以全部子极限组成 $[-1,1]$，上、下极限分别为 $1,-1$。有理数与无理数在每个开区间内都存在，故 $D$ 在任意点的上极限为一，下极限为零；其取值只有零与一，因此子极限只有这两个数，并不包含它们之间的其他数。
\end{solve}

这两个例子说明，子极限集的最大、最小值不能决定中间是否填满；后者还需要函数的其他性质。对连续振荡函数，可以通过适当的趋近点得到整段子极限；对不连续函数，子极限集可以只包含有限个点。

\subsection{估计规则与典型应用}

若 $f\le g$ 在某个充分小的去心邻域内成立，则 $\limsup f\le\limsup g$、$\liminf f\le\liminf g$。局部有界函数的和、非负函数的积满足与式~\eqref{eq:pei-limsup-sum-bounds} 及其乘法版本完全相同的不等式；证明时可直接对同一去心邻域中的确界作估计，或选取逼近边界的数列。若 $h(x)\to c\in\mathbb{R}$，则
\[
    \limsup(f+h)=\limsup f+c,\qquad
    \liminf(f+h)=\liminf f+c.
\]
特别，趋零加性误差不会改变上、下极限。乘法仍需核对极限因子的符号，并在负号出现时交换上、下极限。

\begin{example}
\label{ex:pei-rational-oscillation}
对 $x>1$，令 $f(x)=(x\sin^2x-1)/(x-\sin x)$，求 $x\to+\infty$ 时的上、下极限。
\end{example}
\begin{solve}
先隔离振荡主项：
\[
    f(x)-\sin^2x=\frac{\sin^3x-1}{x-\sin x},\qquad
    |f(x)-\sin^2x|\le\frac2{x-1}\longrightarrow0.
\]
又 $0\le\sin^2x\le1$，在 $x_n=\pi/2+n\pi$ 处取一，在 $y_n=n\pi$ 处取零，所以 $\limsup f=1$、$\liminf f=0$。分母最终离开零点的界保证了误差估计适用于所有足够大的实数，而非只在两条选定点列上成立。
\end{solve}

\begin{example}
\label{ex:pei-dyadic-difference}
设 $f$ 在零点的某个去心邻域中有定义，$F(x)=f(x)/x$ 在那里有界，并且
\[
    \frac{f(x)-f(x/2)}x\longrightarrow A\in\mathbb{R}
    \qquad(x\to0).
\]
证明 $f(x)/x\to2A$。
\end{example}
\begin{solve}
令 $H(x)=[f(x)-f(x/2)]/x$，则恒等式为 $F(x)=H(x)+F(x/2)/2$。记 $U=\limsup_{x\to0}F(x)$、$L=\liminf_{x\to0}F(x)$，有界性保证它们有限；变量替换 $t=x/2$ 不改变零点处的上、下极限。因 $H\to A$，分别取上、下极限得 $U=A+U/2$、$L=A+L/2$，所以 $U=L=2A$，结论成立。这里假设的是商 $f(x)/x$ 有界，不能把它替换成较弱的 $f(x)$ 有界。
\end{solve}

% 第 1.8 节：PDF 第 95--100、102--106 页；第 101 页为夹入的分册目录。
\section{实数及其基本定理}
\label{sec:pei-real-completeness}

\subsection{实数构造与完备性的作用}

有理数具有加、减、乘、除和大小比较，但并不能保证每个有界单调列或每个 Cauchy 列都在有理数中有极限。例如逐位截取 $\sqrt2$ 所得到的有理数近似列满足 Cauchy 条件，却不可能收敛到有理数。引入实数的目标是补足这些由有理近似逼近的位置，同时保留原有运算与次序。

教材介绍了小数、分割和基本列三种表述。用无限小数描述实数时，需要识别 $0.999\ldots=1.000\ldots$ 等不同表示。用有理数的\emph{Dedekind 分割}（Dedekind cut）构造时，以一个非空、真包含于 $\mathbb{Q}$、向下封闭且没有最大元的有理数子集表示一个数；有理数 $r$ 对应 $\{q\in\mathbb{Q}:q<r\}$，其余合适的分割表示无理数。用基本列构造时，则把两个有理 Cauchy 列 $\{q_n\}$、$\{r_n\}$ 在 $|q_n-r_n|$ 最终小于任意给定正有理数时视为等价，并以等价类作为实数。这种定义无需事先假定它们已有实数极限；运算由代表列逐项进行，再证明不依赖代表列的选择。这里给出的是构造思路，完整的域与序公理验证需另行展开。

本章采用公理观点：$\mathbb{R}$ 是满足上确界性质的有序域。其\emph{Archimedes 性质}为：对任意 $x,y>0$，存在 $n\in\mathbb{N}$ 使 $nx>y$。这也能由上确界性质推出：若 $\mathbb{N}$ 有上确界 $s$，则存在整数 $n>s-1$，于是整数 $n+1>s$，矛盾；再对 $y/x$ 选取足够大的整数即可。与有理数相比，新增的实质是\emph{完备性}（completeness），而不是增加另一套形式运算法则。

\subsection{八个基本定理的准确陈述}

下面把本章使用的存在性依据集中列出。为避免把“序列中重复出现的同一个值”与“数集中不同点的积聚”混同，先给出两种对象的定义。

\begin{definition}
\label{def:pei-cluster-point}
若对每个 $\varepsilon>0$ 和每个 $N$，都存在 $n>N$ 使 $|x_n-c|<\varepsilon$，则称 $c$ 为数列 $\{x_n\}$ 的一个\emph{子列极限点}，简称本章的数列极限点；它等价于存在子列 $x_{n_k}\to c$。若对每个 $\varepsilon>0$，数集 $E$ 都在 $0<|x-c|<\varepsilon$ 内有点，则称 $c$ 为 $E$ 的\emph{聚点}（accumulation point）。聚点不要求属于 $E$。若 $c\in E$ 且某个去心邻域内没有 $E$ 的点，则称 $c$ 为\emph{孤立点}（isolated point）。
\end{definition}

数集聚点的定义等价于每个邻域内有无穷多个不同的 $E$ 中点：若只有有限个，可取小于它们到 $c$ 的全部正距离的邻域，排除所有点而产生矛盾。数列的极限点则按指标出现次数计数；例如常数列 $x_n=1$ 的极限点为一，但其值域集合 $\{1\}$ 没有聚点。

\begin{theorem}[确界定理]
\label{thm:pei-supremum}
非空实数集有上界则有实数上确界；有下界则有实数下确界。
\end{theorem}

\begin{theorem}[单调有界定理]
\label{thm:pei-completeness-monotone}
实数列单调不减且有上界，或单调不增且有下界时，必有有限极限。
\end{theorem}

\begin{theorem}[基本列收敛定理]
\label{thm:pei-completeness-cauchy}
每个实数 Cauchy 列都收敛到实数。结合 Cauchy 条件的必要性，这正是定理~\ref{thm:pei-cauchy}。
\end{theorem}

\begin{theorem}[Bolzano--Weierstrass 定理]
\label{thm:pei-bolzano-weierstrass}
每个有界实数列都有收敛子列，其极限仍为实数。这是教材所称的致密性定理。
\end{theorem}

\begin{theorem}[聚点原理]
\label{thm:pei-accumulation}
每个有界无限实数集至少有一个实数聚点。
\end{theorem}

\begin{theorem}[闭区间套定理]
\label{thm:pei-nested-intervals}
设非空有界闭区间 $I_n=[a_n,b_n]$ 满足 $I_{n+1}\subseteq I_n$。则 $\bigcap_nI_n\ne\varnothing$；若还满足 $b_n-a_n\to0$，则公共点唯一。
\end{theorem}

\begin{theorem}[有限覆盖定理]
\label{thm:pei-finite-cover}
设开区间族 $\mathcal{U}$ 覆盖有限闭区间 $[a,b]$，即 $[a,b]\subseteq\bigcup_{U\in\mathcal{U}}U$。则存在有限个 $U_1,\ldots,U_m\in\mathcal{U}$，仍有 $[a,b]\subseteq\bigcup_{j=1}^mU_j$。这些区间称为一个有限子覆盖。
\end{theorem}

\begin{theorem}[Dedekind 分割定理]
\label{thm:pei-dedekind}
设非空集合 $A,B$ 把 $\mathbb{R}$ 分为两部分，$A\cap B=\varnothing$、$A\cup B=\mathbb{R}$，且每个 $a\in A$、$b\in B$ 都满足 $a<b$。则存在唯一实数 $c$ 分隔两组：或者 $c=\max A$ 且 $B$ 无最小值，或者 $c=\min B$ 且 $A$ 无最大值。
\end{theorem}

\emph[caution]{区间套的闭性和长度条件承担不同作用。}没有闭性时，即使长度趋零也可能没有公共点，例如 $(0,1/n)$；有闭性而没有长度趋零时，公共部分可含多个点，例如恒定区间 $[0,1]$。有限覆盖定理同样依赖闭且有界的对象：$(0,1)$ 被 $\{(1/n,1):n\ge2\}$ 覆盖，却不存在有限子覆盖。

\subsection{等价性的主证明链}

这些定理在有序域及 Archimedes 性质的基础上从不同角度表述同一完备性。下面给出一条闭合的证明链，并把聚点、覆盖和分割接入其中；引用前面定理时，只使用当前箭头已经建立的结论。

\begin{proof}[确界、单调列、区间套、子列与基本列之间的推导]
确界定理推出单调有界定理，证明已在定理~\ref{thm:pei-monotone} 中给出。单调有界定理推出区间套定理：$a_n$ 递增有上界、$b_n$ 递减有下界，故分别趋于 $a,b$，且对任意 $n,m$ 有 $a_n\le b_m$，所以 $a\le b$。逐项取极限可知公共部分恰为 $[a,b]$；若长度趋零，则 $a=b$。

区间套定理推出有界列的收敛子列：取一个含所有项的闭区间，将其二等分，至少一半含有无穷多个指标对应的项；保留该半区间并反复进行，得到长度趋零的闭区间套 $I_k$。设唯一公共点为 $c$，每次选择比前次更大的指标 $n_k$，使 $x_{n_k}\in I_k$，便有 $|x_{n_k}-c|\le|I_k|\to0$。这里“无穷多项”按指标计数，允许项值重复。

有界列的收敛子列性质推出 Cauchy 列收敛：基本列先必有界，因为在误差取一后，其尾部落在某个固定项的一邻域内，再并入有限首段即可。取收敛子列 $x_{n_k}\to c$；给定 $\varepsilon>0$，用 Cauchy 条件控制两项之差小于 $\varepsilon/2$，再选一个充分后期且距离 $c$ 小于 $\varepsilon/2$ 的子列项，三角不等式便使全部尾项距离 $c$ 小于 $\varepsilon$。

最后仅用基本列收敛性质反推确界定理。设 $E\ne\varnothing$ 且有上界 $b_0$；若某个 $a_0\in E$ 已是上界，它就是最大值及上确界。否则 $a_0$ 不是上界，$a_0<b_0$。逐次二分 $[a_n,b_n]$：中点若为上界，就替换右端；若不是上界，就替换左端。始终保持 $a_n$ 不是上界、$b_n$ 是上界，且长度为 $2^{-n}(b_0-a_0)$。两端点列都是 Cauchy 列，故由所假定的基本列收敛性质有共同极限 $c$。若某个 $x\in E$ 大于 $c$，则充分后期有 $b_n<x$，与上界性质矛盾；若 $y<c$，则充分后期有 $a_n>y$，而 $a_n$ 不是上界，故存在 $x\in E$ 使 $x>a_n>y$。因此 $c$ 是最小上界，闭合证明链完成。
\end{proof}

聚点原理与收敛子列性质的等价需要注意重复值。由有界无限集选出两两不同的点列，再选其收敛子列，极限即为原集聚点。反过来，若有界数列的某个值出现无穷多次，直接取常值子列；否则值域必为无限有界集。取其聚点 $c$，每个邻域有无穷多个不同值，因而有任意大的原指标，可递归选出趋于 $c$ 的子列。这两种情况涵盖全部有界数列。

\begin{proof}[区间套推出有限覆盖，有限覆盖推出收敛子列]
若某个覆盖 $[a,b]$ 的开区间族没有有限子覆盖，将 $[a,b]$ 二等分，至少有一半仍没有有限子覆盖；不断保留这样的半区间，得到长度趋零的闭区间套。其公共点 $c$ 属于某个原覆盖区间 $U$，而 $U$ 是开集，故存在 $r>0$ 使 $(c-r,c+r)\subseteq U$。充分后期的套区间长度小于 $r$，又包含 $c$，所以整个套区间都在 $U$ 内，竟由单个区间覆盖，矛盾。

反向，设有界列 $x_n\in[a,b]$ 没有收敛子列。对每个 $c\in[a,b]$，必存在一个开邻域 $U_c$，其中只含有限多个指标对应的项；否则按半径 $1/k$ 逐次选取任意大指标，就能得到趋于 $c$ 的子列。这些 $U_c$ 覆盖 $[a,b]$，有限覆盖定理给出有限子覆盖，因而全部指标只能落在有限个有限指标集中，矛盾。于是有限覆盖性质也接入了上述等价链。
\end{proof}

\begin{proof}[确界与分割之间的推导]
给定定理~\ref{thm:pei-dedekind} 的分割 $A,B$，集合 $A$ 非空并被 $B$ 的任意元素界住，令 $c=\sup A$。每个 $a\in A$ 都不超过 $c$，每个 $b\in B$ 都不小于 $c$。由于 $A,B$ 分割全部实数，$c$ 必属于其中一组；属于 $A$ 时是其最大值，属于 $B$ 时是其最小值。两组不能同时有最大值和最小值，因为若两者不同，它们的中点将无处归属；若相同，则违反两组不交。分割点的唯一性也由同一中点论证得到。

反向，设非空集 $E$ 有上界，令 $B$ 为全部上界组成的集合，$A$ 为全部非上界组成的集合。两组均非空，且任意 $a\in A$ 之后存在 $x\in E$，从而 $a<x\le b$ 对每个 $b\in B$ 成立，所以它们构成分割。$A$ 不可能有最大值：对 $a\in A$，取 $x\in E$ 且 $x>a$，则 $(a+x)/2$ 仍不是上界且大于 $a$。分割定理因而保证 $B$ 有最小值，它正是 $E$ 的上确界。
\end{proof}

\subsection{局部性质、整体结论与极限点集}

\begin{example}
\label{ex:pei-local-boundedness-compact}
若函数 $f:[a,b]\to\mathbb{R}$ 在每点附近相对于 $[a,b]$ 局部有界，证明它在整个区间上有界。
\end{example}
\begin{solve}
对每个 $x\in[a,b]$，取开邻域 $U_x$ 与有限常数 $M_x$，使 $t\in U_x\cap[a,b]$ 时 $|f(t)|\le M_x$。这些邻域构成开覆盖，取有限子覆盖 $U_{x_1},\ldots,U_{x_m}$，便有 $|f(t)|\le\max_jM_{x_j}$ 对全区间成立。因此每点存在有限双侧极限、端点存在相应有限单侧极限的函数，即使不连续，也在 $[a,b]$ 上有界：极限控制去心邻域，再把该点有限的函数值并入上界即可。若区间改为 $(0,1)$，$f(t)=1/t$ 在每点局部有界却整体无界，说明闭区间条件的作用。
\end{solve}

\begin{example}
\label{ex:pei-bisection-fixed-point}
设 $f:[0,1]\to\mathbb{R}$ 连续，且 $f(0)>0$、$f(1)<1$。用区间套构造 $c\in(0,1)$ 使 $f(c)=c$。
\end{example}
\begin{solve}
令 $g(x)=f(x)-x$，则 $g(0)>0$、$g(1)<0$。二分 $[0,1]$，中点若使 $g=0$，结论已得；否则保留左端为正、右端为负的半区间。反复进行得到闭区间套 $[a_n,b_n]$，长度趋零且 $g(a_n)>0$、$g(b_n)<0$。其唯一公共点为 $c$，连续性给出 $g(c)\ge0$ 及 $g(c)\le0$，所以 $g(c)=0$；端点处符号严格，故 $c\ne0,1$。本例用连续性把区间内的逼近转化为方程的精确解，与例~\ref{ex:pei-monotone-fixed-point} 的单调映射证明使用了不同假设。
\end{solve}

\begin{proposition}
\label{prop:pei-closed-cluster-set}
设 $C$ 是一个实数列全部有限子列极限点组成的集合，则 $C$ 是闭集；若原数列有界，$C$ 还是非空有界闭集，且 $C$ 只有一个元素当且仅当原数列收敛。
\end{proposition}
\begin{proof}
设 $c_j\in C$ 且 $c_j\to c\in\mathbb{R}$。因每个 $c_j$ 都是原数列的子列极限，可递归选取 $n_j>n_{j-1}$，使 $|x_{n_j}-c_j|<1/j$。于是 $|x_{n_j}-c|\le1/j+|c_j-c|\to0$，故 $c\in C$，证明了闭性。有界性由原列传给各子列极限，非空性由 Bolzano--Weierstrass 定理保证。

若有界原列仅有一个子列极限点 $c$ 而不趋于 $c$，则存在 $\varepsilon_0>0$ 及子列使 $|x_{n_j}-c|\ge\varepsilon_0$。该子列仍有界，故存在收敛子列，其极限 $d$ 满足 $|d-c|\ge\varepsilon_0$，与唯一性矛盾。反向由收敛列的全部子列都具有相同极限得到。
\end{proof}

例如，不存在有限子列极限点恰好为 $\{1,1/2,1/3,\ldots\}$ 的实数列，因为该集合不闭，漏掉了零。另一方面，“唯一有限子列极限点”若没有有界性并不足以保证收敛：$x_{2n}=0$、$x_{2n-1}=n$ 的唯一有限子列极限点为零，原列仍不收敛。与第~\ref{sec:pei-sequence-limsup} 节允许无穷子列极限的表述相对照，可以看出有界性正是在排除逃向无穷大的可能。


% 教材来源：数学分析/textbooks/数分习题（分章原文）/数分习题-2.pdf。
% PDF 第 1--19 页对应第 2.1 节，第 19--32 页对应第 2.2 节，
% 第 32--33 页对应第 2.3 节，第 33--42 页对应第 2.4 节。




